\section{Cavity resonances: the generalised eigenproblem and its null space} \label{sec:eigen}

\noindent\emph{Status: draft -- the eigenproblem, the free-edge restriction and the dense solve are
implemented and covered by tests; the computed spectrum of a rectangular cavity lies below the
analytic one by $0.6$ per cent on the fine grid and the gradient null space is removed by a
static-eigenvalue cut. The deflation that would let the same problem be solved iteratively is not
implemented, so a grid whose free-edge count exceeds \texttt{max\_dense} is refused rather than
approximated.}

The two Maxwell grid equations \eqref{eq:faraday} and \eqref{eq:ampere} of Section~\ref{sec:theory}
are a generalised eigenproblem in disguise: assume harmonic time dependence, drop the source,
eliminate one of the two fields, and what remains is a statement about the primal curl and the two
Hodge stars alone. Everything that follows turns on the one structural fact proved in
Proposition~\ref{prop:identities}: the curl annihilates gradients. On the full complex that is a
property one is glad to have. After the PEC restriction discussed in Section~\ref{sec:theory} it
becomes a nuisance, because the restricted curl--curl operator inherits a null space that is neither
empty nor small -- 216 of the 762 unknowns on the coarse cavity grid. No spectrum can be believed
before those modes are identified. Section~\ref{sec:formulation} states the eigenproblem and the
symmetrisation the code actually solves; Section~\ref{sec:nullspace} characterises the null space
and settles which test for it works; Section~\ref{sec:spectrum} compares the surviving frequencies
with the analytic spectrum of a conducting box and measures their convergence;
Section~\ref{sec:rejected} records the two approaches that failed, because a failure with a reason
is worth more than a heuristic that happens to work.

\subsection{Formulation} \label{sec:formulation}

With the phasor convention $\mathbf{e}(t)=\hat{\mathbf{e}}\exp(-\mathrm{i}\omega t)$ -- the only
place in this
report where complex fields are used -- and with no source, $\mathbf{j}=\mathbf{0}$, the grid
equations
\eqref{eq:faraday} and \eqref{eq:ampere} become a pair of algebraic relations between phasors,
\begin{equation}
  \Cmat\hat{\mathbf{e}} = \mathrm{i}\omega\Mm\hat{\mathbf{h}}, \qquad
  \Cmat^{\!\top}\hat{\mathbf{h}} = -\mathrm{i}\omega\Em\hat{\mathbf{e}}.
  \label{eq:harmonic}
\end{equation}
Eliminating $\hat{\mathbf{h}}$ between the two lines of \eqref{eq:harmonic} leaves a statement
about the electric
field alone,
\begin{equation}
  \Cmat^{\!\top}\Mm^{-1}\Cmat\,\hat{\mathbf{e}} = \omega^{2}\Em\,\hat{\mathbf{e}},
  \label{eq:geneig}
\end{equation}
a generalised symmetric eigenproblem with the same curl--curl operator whose time-domain
counterpart we met in
Section~\ref{sec:energy}. The two matrices on the left and right of \eqref{eq:geneig} are symmetric
and
$\Em$ is positive definite because both stars are diagonal with positive entries \eqref{eq:stars},
so the pencil
is self-adjoint in the $\Em$ inner product. Two consequences are used below without further
comment: the
eigenvalues $\omega^{2}$ are real and non-negative, and eigenvectors belonging to distinct
eigenvalues are
$\Em$-orthogonal. Eliminating $\hat{\mathbf{e}}$ instead gives the dual statement with $\Mm$ and
$\tilde\Cmat$ in the roles of $\Em$ and $\Cmat$; the two formulations have the same spectrum, and
we keep the
electric one because $\Em$ is diagonal and cheap to invert.

The conductor condition \eqref{eq:pec} pins the tangential voltage on every wall edge to zero,
which is a
constraint on the unknown space rather than on the equation. Writing $\mathcal{F}$ for the set of
free edges and
$\Cmat_f=\Cmat_{:,\mathcal{F}}$ for the curl with the pinned columns deleted, \eqref{eq:geneig}
restricted to the
unknowns reads
\begin{equation}
  \Cmat_f^{\!\top}\Mm^{-1}\Cmat_f\,\hat{\mathbf{e}}_f = \omega^{2}\Em_f\,\hat{\mathbf{e}}_f,
  \label{eq:geneig-f}
\end{equation}
with $\Em_f$ the corresponding submatrix of the diagonal star. Only edges are restricted: a perfect
conductor
constrains the primal edge space, whereas a perfect magnetic conductor would constrain the primal
face space, so
$\Mm$ and the face count are untouched. This is the same asymmetry that the two Gauss laws exhibit,
and it is the
reason the free-edge restriction costs us a null space while the dual restriction would not.

Equation \eqref{eq:geneig-f} is a symmetric pencil but not a standard problem, and the dense
routine we call
expects the latter. Since $\Em_f$ is diagonal and positive we multiply through by $\Em_f^{-1/2}$,
\begin{equation}
  \Em_f^{-1/2}\Cmat_f^{\!\top}\Mm^{-1}\Cmat_f\Em_f^{-1/2}\,\mathbf{y} = \omega^{2}\mathbf{y},
  \qquad \hat{\mathbf{e}}_f=\Em_f^{-1/2}\mathbf{y},
  \label{eq:symm}
\end{equation}
and are left with a symmetric positive semi-definite matrix whose eigenvalues are the squared
angular frequencies
directly and whose eigenvectors are the scaled edge voltages. The semi-definiteness is not a defect
of the
scaling: the matrix in \eqref{eq:symm} is a Gram matrix of $\Mm^{-1/2}\Cmat_f\Em_f^{-1/2}$ and its
kernel is
exactly the null space of Section~\ref{sec:nullspace}.

The code assembles the sparse product $\Cmat^{\!\top}\Mm^{-1}\Cmat$ once, takes its free--free
block, and hands it to \texttt{scipy.linalg.eigh}, which returns the complete spectrum. We chose a
dense solver deliberately, and the reasons are worth stating because they will have to be given up
one at a time. It is exact in the sense that it needs no convergence argument, no shift, no
preconditioner and no estimate of how many eigenvalues the caller wants: there is no way for it to
return a plausible wrong answer. It also returns the null cluster, which we want, because its
dimension is the quantity we check against $(n_x-1)(n_y-1)(n_z-1)$. The cost is cubic, and at the
sizes that matter here it is irrelevant: the largest grid in this report has 2355 unknowns, a
symmetric eigenproblem of that order costs of order $2355^{3}\approx1.3\times10^{10}$ floating-point
operations, and one solve takes seconds. Above \texttt{max\_dense} $=3000$ the routine raises
instead of trying, which is the behaviour we want from a validation instrument.
Table~\ref{tab:sizes} lists the sizes.

\begin{table}[H]
  \centering
  \caption{Size of the restricted eigenproblem \eqref{eq:symm} for the two cavity grids used in
this report, and
  the split between the gradient null space and the physical unknowns. The last two columns sum to
the third.}
  \label{tab:sizes}
  \begin{tabular}{@{}lrrrr@{}}
    \toprule
    grid & primal edges & free edges & gradient null modes & physical unknowns \\
    \midrule
    $10\times7\times5$ & 1382 & 762  & 216 & 546  \\
    $14\times10\times7$ & 3587 & 2355 & 702 & 1653 \\
    \bottomrule
  \end{tabular}
\end{table}

\subsection{The gradient null space} \label{sec:nullspace}

A gradient field $\Gmat\boldsymbol\varphi$ is annihilated by the curl, $\Cmat\Gmat=\mathbf{0}$, so
every such
field is a solution of \eqref{eq:geneig} with $\omega=0$. The question is which of them survive the
restriction.
A vector in the restricted unknown space must vanish on the pinned edges, and a gradient field
$\Gmat\boldsymbol\varphi$
does so exactly when the potential $\varphi$ is equal at the two endpoints of every wall edge. The
pinned edges
connect the boundary nodes to each other, so this forces $\varphi$ to be constant on the boundary,
and a constant
is no field at all; what remains is a potential that vanishes on every boundary node and is free at
the interior
nodes. There are exactly $(n_x-1)(n_y-1)(n_z-1)$ of those, and each gives an independent gradient
field, so
\begin{equation}
  \dim\ker\Cmat_f = (n_x-1)(n_y-1)(n_z-1) = 9\cdot6\cdot4 = 216
  \label{eq:nullcount}
\end{equation}
on the coarse grid, the count of \eqref{eq:nullcount}, and $13\cdot9\cdot6=702$ on the fine one.
The solver reports
exactly those counts as rejected static modes, on every grid we have tried.

It is tempting to conclude from the preceding paragraph that the restriction is what removes the
null space. It is not, and the error is instructive because it is easy to make in a code review. The
identity $\Cmat\Gmat=\mathbf{0}$ is an identity on the \emph{full} complex, and the proof in
Proposition~\ref{prop:identities} is a telescoping sum around the loop bounding a face. Deleting the
columns of the pinned edges deletes precisely those terms of the sum that cancel, so the restricted
product $\Cmat_f\Gmat_f$ is not zero, and there is no restricted identity to appeal to. On the full
1382-edge space there are $528-1=527$ independent gradients; after the restriction 216 of them are
still in the kernel, and the remaining 311 are not in the unknown space at all, since their pinned
entries are non-zero. The restriction shrinks the null space; it does not remove it.

The second test we tried, and the second one to be discarded, was the divergence residual. Since a
physical mode of the ideal problem satisfies $\Gmat_f^{\!\top}\Em_f\hat{\mathbf{e}}_f=\mathbf{0}$,
one may hope that $\|\Gmat_f^{\!\top}\Em_f\hat{\mathbf{e}}_f\|/\|\Em_f\hat{\mathbf{e}}_f\|$
separates the two families. Measured on the coarse grid, the gradients return $2.40$ to $2.68$ and
the physical modes $0.45$ to $0.78$: two clusters of the same order of magnitude, with no
machine-precision zero anywhere. The reason is structural rather than numerical. The only discrete
identity available is $\Smat\Cmat=\mathbf{0}$ in \eqref{eq:identities}, and it relates the
divergence to the \emph{curl}, not to the gradient; a discrete analogue of
$\nabla\cdot(\eps\nabla\varphi)=0$ would require $\Smat\Em\Gmat=\mathbf{0}$, which fails because
$\Em$ is a metric and not the identity. A gradient therefore has no reason to be distinguished by
any divergence test, and none is. The implementation still applies the residual as a screen with a
threshold of $1.5$, between the two clusters, but it is a belt over braces: on every grid of this
report it rejects nothing, and the tests assert that it rejects nothing.

What does work is the crudest test available. The null cluster sits at rounding level, and the
physical spectrum
does not, so a threshold on the eigenvalue separates them,
\begin{equation}
  \text{accept } \omega^{2} > \tau\,\omega^{2}_{\mathrm{ref}}, \qquad \tau=10^{-10}, \qquad
  \omega^{2}_{\mathrm{ref}} = \left(\frac{\pi c_0}{L_{\max}}\right)^{2},
  \label{eq:cut}
\end{equation}
with $L_{\max}$ the largest box dimension, $1.0$ m. On the coarse grid the reference frequency is
$c_0/2L_{\max}=149.9$ MHz; the null cluster is measured at
$1.3\times10^{-13}\,\omega^{2}_{\mathrm{ref}}$ and the
fundamental at $3.0\,\omega^{2}_{\mathrm{ref}}$. The gap is therefore thirteen orders of magnitude
wide, and the
threshold sits about three orders above the cluster and ten below the fundamental -- it could be
moved by an
order in either direction without changing a single accepted mode. This is the most load-bearing
number in the
solver, and it is the one we would least like to lose: $\tau$ is a statement about floating-point
arithmetic, not
about the physics, and on a strongly graded mesh, where the assembly error may spread the cluster,
it will need
to become a relative gap test between consecutive eigenvalues.

\subsection{The spectrum of a conducting box} \label{sec:spectrum}

A rectangular cavity with perfectly conducting walls has the textbook spectrum
\begin{equation}
  f_{mnp} = \frac{c_0}{2}\sqrt{\left(\frac{m}{a}\right)^{2}+\left(\frac{n}{b}\right)^{2}
  +\left(\frac{p}{d}\right)^{2}},
  \label{eq:analytic}
\end{equation}
for a box $a\times b\times d=1.0\times0.7\times0.5$ m, and it is our reference because it exercises
the
discretisation on three different length scales at once. The counting is the part that is usually
got wrong.
Triples with a single non-zero index are not Maxwell modes and must be dropped: a TM mode needs
$m,n\ge1$ and a
TE mode needs $p\ge1$ together with $m\ge1$ or $n\ge1$. A triple with exactly two non-zero indices
is one mode --
$(m,n,0)$ is TM, $(m,0,p)$ and $(0,n,p)$ are TE -- and a triple with all three non-zero is TE/TM
degenerate and
counts twice. The lowest six modes are then $261.388$, $335.178$, $368.416$ twice and $397.743$
twice, in MHz, and
the doubling has two different causes that we keep apart. The pair at $397.743$ MHz is the genuine
degeneracy of
$(1,1,1)$, one TE and one TM mode that any rectangular box has. The pair at $368.416$ MHz is an
accident of the
proportions: because $b=d=a/2$, the triples $(0,1,1)$ and $(2,1,0)$ give the same value
$\sqrt{6.0408}\,c_0/2a$ even though they are one TE and one TM mode of different shape. A box with
$b\neq d$
would separate them. Counting the accident as a degeneracy is right for this box and wrong as a
rule, and the
reference list intentionally contains the duplicate so that it can be compared one-to-one with a
computed
spectrum that resolves both polarisations.

\begin{table}[H]
  \centering
  \caption{Lowest six resonances of the $1.0\times0.7\times0.5$ m cavity: analytic values from
\eqref{eq:analytic}
  and the computed ones on the two grids. Errors are relative and in per cent. Every computed value
lies below
  its analytic counterpart, as the variational argument requires.}
  \label{tab:cavity}
  \begin{tabular}{@{}lrrrrrr@{}}
    \toprule
    mode & analytic [MHz] & mult.\ & coarse [MHz] & err [\%] & fine [MHz] & err [\%] \\
    \midrule
    $(1,1,0)$          & 261.388 & 1 & 259.567 & 0.697 & 260.487 & 0.345 \\
    $(1,0,1)$          & 335.178 & 1 & 330.518 & 1.390 & 332.794 & 0.711 \\
    $(0,1,1)$, $(2,1,0)$ & 368.416 & 2 & 363.384 & 1.366 & 365.863 & 0.693 \\
    $(1,1,1)$          & 397.743 & 2 & 392.852 & 1.230 & 395.260 & 0.624 \\
    \bottomrule
  \end{tabular}
\end{table}

Table~\ref{tab:cavity} puts the two spectra side by side. Two features matter more than the
individual numbers. First, every computed frequency is \emph{below} its analytic counterpart,
without exception. That is not luck: the discretisation is a Rayleigh--Ritz method on a subspace of
the continuous solution space, so its eigenvalues bound the exact ones from above in $\omega^{2}$
and the frequencies from below \cite{bossavit1998,jin2014}. A method that overshot the analytic
value on a converged grid would be a method with a sign error somewhere, and the bound gives us a
check that costs nothing. Second, the errors are of order one per cent on the coarse grid, which is
what the dispersion analysis of Section~\ref{sec:dispersion} leads one to expect: the fundamental is
resolved by about ten cells per wavelength along the long side, and the deficit is second order in
$h/\lambda$.

The comparison in Table~\ref{tab:cavity} covers the six lowest modes because they are the ones the
reference list resolves unambiguously. Beyond them the matching becomes a question about resolution
rather than about the operator: the seventh and eighth analytic values, $423.971$ and $453.749$ MHz,
come out at $417.031$ and $440.131$ MHz on the coarse grid -- $1.6$ and $3.0$ per cent -- and at
$420.421$ and $447.033$ MHz on the fine one, because $(1,2,0)$ resolves only three and a half cells
per half wavelength along the $0.7$ m side. Those errors belong to the same second-order family, and
we quote them rather than hide them, but they are not part of the mean below.

\subsection{Convergence under refinement}

The refinement between the two grids is a factor $1.4$ in every direction: $h=0.1$ m on the coarse
grid and $h=0.1/1.4=0.0714$ m on the fine one, with cells that are cubic to within two per cent, so
the star metric \eqref{eq:stars} stays isotropic and the comparison isolates the mesh size. The mean
relative error over the six modes falls from $1.213$ to $0.615$ per cent, a ratio of $0.507$, and
the worst-case error from $1.39$ to $0.711$ per cent. The ratio predicted by a second-order method
is $1/1.4^{2}=0.510$. Table~\ref{tab:convergence} collects the numbers.

\begin{table}[H]
  \centering
  \caption{Convergence of the six lowest resonances of Table~\ref{tab:cavity} under uniform
refinement by a
  factor $1.4$. The error ratio is compared with the $h^{2}$ prediction.}
  \label{tab:convergence}
  \begin{tabular}{@{}lrrrr@{}}
    \toprule
    grid & $h$ [m] & free edges & mean err [\%] & worst err [\%] \\
    \midrule
    $10\times7\times5$   & 0.1000 & 762  & 1.213 & 1.390 \\
    $14\times10\times7$  & 0.0714 & 2355 & 0.615 & 0.711 \\
    \midrule
    ratio                & 0.714  &      & 0.507 & 0.512 \\
    prediction $h^{2}$   &        &      & 0.510 & 0.510 \\
    \bottomrule
  \end{tabular}
\end{table}

The agreement between $0.507$ and $0.510$ is the strongest evidence in this report that the
discretisation is right rather than tuned. A fitted correction factor, an empirically adjusted star,
or a reference spectrum adjusted to fit the output would all reproduce the coarse numbers; none of
them would halve the error when the cell size is divided by $1.4$. Second order is also what the
construction promises -- the only approximation in the method is the diagonal star, and the FIT
operators themselves are exact \cite{weiland1977,clemens2001} -- so the measurement confirms the
analysis rather than standing in for it. The caveat is that both grids are uniform, and second order
has not been demonstrated on a graded mesh; that remains an open question.

\subsection{What we tried and rejected} \label{sec:rejected}

Everything above is a dense solve. A dense solve is not available on the grids we want next, and
the obvious
alternative -- an iterative Krylov eigensolver, ARPACK's \texttt{eigsh} with shift-invert -- fails
here in three
distinct ways, each of which produces a plausible-looking answer rather than an error:
\begin{itemize}[leftmargin=1.4em,itemsep=2pt,topsep=2pt]
  \item With the shift at $\sigma=0$ the iteration converges onto the null space, since the null
space sits
        exactly at the shift and is 216- to 702-dimensional; a request for six vectors returns
mostly static
        modes, and too few physical ones survive to be called a spectrum.
  \item Shifting upwards to escape the null space is worse, not better: shift-invert finds the
eigenvalues
        nearest the shift, so a shift placed above the fundamental returns a window \emph{above} it
and the
        lowest mode is silently missing. The output looks like a cavity spectrum with the
fundamental absent,
        which is exactly the failure mode hardest to notice.
  \item Escalating or overlapping several shifts does not repair this, because no choice of shift
makes a Krylov
        method find eigenvalues below its target. The difficulty is not the shift but the order in
which the
        method visits the spectrum, and the null space is always visited first.
\end{itemize}
The conclusion is that the null space has to be removed \emph{before} the iteration, not filtered
afterwards and
not avoided by shifting. The fix is a deflation: project the gradient subspace out with the
$\Em$-orthogonal
projector built from a sparse potential solve,
$\mathbf{P}=\mathbf{I}-\Gmat_f(\Gmat_f^{\!\top}\Em_f\Gmat_f)^{-1}
\Gmat_f^{\!\top}\Em_f$, and run a block method such as LOBPCG on the deflated operator. The
operator in brackets
is a discrete Laplacian on the interior nodes, so applying the projector costs one sparse solve per
iteration --
cheap, but not yet written. That work is milestone M5 of the roadmap; until it lands, dense
\texttt{eigh} ships
with the \texttt{max\_dense} guard, which covers every grid in this report and none of the ones we
want next.

One further idea deserves a paragraph, because it looks like a filter and is not one. Define the
curl residual
\begin{equation}
  r(\hat{\mathbf{e}}) = \frac{\|\Cmat\hat{\mathbf{e}}\|}{\|\hat{\mathbf{e}}\|\;\omega/c_0} \ \
[\mathrm{m}],
  \label{eq:rcurl}
\end{equation}
which has the dimensions of a length and was meant to answer the question "is this vector a wave or
junk?" It
does not answer it. The curl $\Cmat$ is dimensionless, so the numerator measures how much the field
changes from
one edge to the next, which for any resolved mode is of order $h/\lambda$ relative to the field
itself; dividing
by $\omega/c_0=2\pi/\lambda$ leaves the cell size and nothing else. Measured on the coarse grid,
\eqref{eq:rcurl}
returns $0.1000$ m -- exactly the pitch $h=0.1$ m -- for the fundamental, for the degenerate pair,
and for every
other physical mode, whatever its frequency or polarisation. It is a rewriting of the eigen-relation
\eqref{eq:geneig}, which pins
$\|\Cmat\hat{\mathbf{e}}\|^{2}_{\Mm^{-1}}=\omega^{2}\|\hat{\mathbf{e}}\|^{2}_{\Em}$
for each eigenvector, so it cannot distinguish one eigenvector from another. The only vectors for
which
\eqref{eq:rcurl} returns something else are the exact gradients, where
$\Cmat\hat{\mathbf{e}}=\mathbf{0}$; those
are precisely the modes the static cut \eqref{eq:cut} already removes at machine precision, so the
residual adds
nothing even there, and for a contaminated vector -- the case it was invented for -- the physical
admixture
dominates the norm and the quotient returns the pitch again. We record it here so that it is not
reinvented.

What we can claim, then, is narrow and we state it narrowly. The FIT eigenproblem is symmetric and
its spectrum is real and bounded below by zero; on a rectangular box the computed frequencies
converge at second order to the analytic ones, from below, with errors of half a per cent on a grid
of 2355 unknowns; the gradient null space is of dimension $(n_x-1)(n_y-1)(n_z-1)$, it is a
consequence of the PEC restriction rather than of the discretisation, and it is separated from the
physical spectrum by more than ten orders of magnitude. What we cannot claim is that any of this
survives an iterative solver, a graded mesh, or a shape that is not a box, and the next section
describes the implementation that the dense path rests on.
